\documentclass[a4paper]{article}

\usepackage[margin=1in]{geometry}
\usepackage{amsmath,amsfonts,mathtools,amsthm,amssymb}
\usepackage[l2tabu,orthodox]{nag}
\usepackage{microtype} % fixes spacing or whatever
\usepackage{enumitem}
\usepackage{tikz-cd}
\usepackage{xcolor}
\usepackage{physics}
\usepackage{mathrsfs}
\usepackage{cancel}
% \usepackage{libertine,libertinust1math}
\usepackage{eucal}
\usepackage[toc,page]{appendix}
\usepackage{pgfplots}
\pgfplotsset{compat=1.18}

\newtheorem{proposition}{Proposition}
\newtheorem{theorem}{Theorem}
\newtheorem{lemma}{Lemma}
\newtheorem{corollary}{Corollary}

\theoremstyle{definition}
\newtheorem{remark}{Remark}
\newtheorem{definition}{Definition}
\newtheorem{example}{Example}

% \usepackage[sc,osf]{mathpazo}   % With old-style figures and real smallcaps.
% \linespread{1.025}              % Palatino leads a little more leading
% % Euler for math and numbers
% \usepackage[euler-digits,small]{eulervm}
% \AtBeginDocument{\renewcommand{\hbar}{\hslash}}

\usepackage{etoolbox}
\usepackage{dynkin-diagrams}
\def\row#1/#2!{\dynkin{#1}{#2}\\}
\newcommand{\tble}[1]{
   \renewcommand*\do[1]{\row##1!}
   \[
      \begin{array}{ll}\docsvlist{#1}\end{array}
   \]
}

\allowdisplaybreaks

\renewcommand{\epsilon}{\varepsilon}
% \renewcommand{\phi}{\varphi}

\DeclareMathAlphabet{\mathpzc}{OT1}{pzc}{m}{it}

\newcommand{\goth}[1]{\mathfrak{#1}}
\newcommand{\red}[1]{#1_{\text{red}}}
\newcommand{\ad}[1]{#1_{\text{Ad}}}
\newcommand{\reg}[1]{#1_{\text{reg}}}
\newcommand{\sh}[1]{#1^{\text{sh}}}
\newcommand{\cpdv}[2]{\cfrac{\partial #1}{\partial #2}}

\newcommand{\fA}{\mathcal{A}}
\newcommand{\fB}{\mathcal{B}}
\newcommand{\fC}{\mathcal{C}}
\newcommand{\fD}{\mathcal{D}}
\newcommand{\fE}{\mathcal{E}}
\newcommand{\fF}{\mathcal{F}}
\newcommand{\fG}{\mathcal{G}}
\newcommand{\fH}{\mathcal{H}}
\newcommand{\fI}{\mathcal{I}}
\newcommand{\fJ}{\mathcal{J}}
\newcommand{\fK}{\mathcal{K}}
\newcommand{\fL}{\mathcal{L}}
\newcommand{\fM}{\mathcal{M}}
\newcommand{\fN}{\mathcal{N}}
\newcommand{\fO}{\mathcal{O}}
\newcommand{\fP}{\mathcal{P}}
\newcommand{\fQ}{\mathcal{Q}}
\newcommand{\fR}{\mathcal{R}}
\newcommand{\fS}{\mathcal{S}}
\newcommand{\fT}{\mathcal{T}}
\newcommand{\fX}{\mathcal{X}}
\newcommand{\fY}{\mathcal{Y}}
\newcommand{\fZ}{\mathcal{Z}}

\newcommand{\be}{\mathbf{e}}
\newcommand{\bx}{\mathbf{x}}
\newcommand{\bone}{\mathbf{1}}
\newcommand{\cx}{\bar{x}}
\newcommand{\cy}{\bar{y}}
\newcommand{\vx}{\vec{x}}
\newcommand{\vsigma}{\vec{\sigma}}
\newcommand{\tzeta}{\tilde{\zeta}}

\newcommand{\gm}{\goth{m}}
\newcommand{\gp}{\goth{p}}
\newcommand{\gF}{\goth{F}}
\newcommand{\gU}{\goth{U}}
\newcommand{\gV}{\goth{V}}

\newcommand{\A}{\mathbf{A}}
\newcommand{\C}{\mathbf{C}}
\newcommand{\F}{\mathbf{F}}
\newcommand{\G}{\mathbf{G}}
\newcommand{\bH}{\mathbf{H}}
\newcommand{\N}{\mathbf{N}}
\newcommand{\PP}{\mathbf{P}}
\newcommand{\Q}{\mathbf{Q}}
\newcommand{\R}{\mathbf{R}}
\newcommand{\Z}{\mathbf{Z}}

\newcommand{\dlog}{\dd\log}

\newcommand\srestr[2]{{\left.\kern-\nulldelimiterspace #1\vphantom{\small|} \right|_{#2}}}
\newcommand\restr[2]{{\left.\kern-\nulldelimiterspace #1 \vphantom{\big|} \right|_{#2}}}
\newcommand\gsec[2]{\Gamma({#1}, {#2})}
\newcommand\gsecx[1]{\Gamma(X, {#1})}

\DeclareMathOperator{\chr}{char}
\DeclareMathOperator{\rProj}{\mathbf{Proj}}
\DeclareMathOperator{\rSpec}{\mathbf{Spec}}
\DeclareMathOperator{\rHom}{\mathbf{Hom}}
\DeclareMathOperator{\rExt}{\mathbf{Ext}}
\DeclareMathOperator{\id}{id}
\DeclareMathOperator{\Frac}{Frac}
\DeclareMathOperator{\rk}{rank}
\DeclareMathOperator{\deter}{det}
\DeclareMathOperator{\pic}{Pic}
\DeclareMathOperator{\cacl}{CaCl}
\DeclareMathOperator{\trd}{tr.d.}
\DeclareMathOperator{\cl}{Cl}
\DeclareMathOperator{\depth}{depth}
\DeclareMathOperator{\codim}{codim}
\DeclareMathOperator{\Div}{Div}
\DeclareMathOperator{\coker}{coker}
\DeclareMathOperator{\len}{length}
\DeclareMathOperator{\height}{height}
\DeclareMathOperator{\supp}{Supp}
\DeclareMathOperator{\proj}{Proj}
\DeclareMathOperator{\im}{im}
\DeclareMathOperator{\Hom}{Hom}
\DeclareMathOperator{\Der}{Der}
\DeclareMathOperator{\spec}{Spec}
\DeclareMathOperator{\Aut}{Aut}
\DeclareMathOperator{\ch}{char}
\DeclareMathOperator{\tor}{Tor}
\DeclareMathOperator{\Ann}{Ann}
\DeclareMathOperator{\Syl}{Syl}
\DeclareMathOperator{\Sym}{Sym}
\DeclareMathOperator{\GL}{GL}
\DeclareMathOperator{\SL}{SL}
\DeclareMathOperator{\Stab}{Stab}
\DeclareMathOperator{\Perm}{Perm}
\DeclareMathOperator{\Orb}{Orb}
\DeclareMathOperator{\Gal}{Gal}
\DeclareMathOperator{\Supp}{Supp}
\DeclareMathOperator{\Ext}{Ext}
\DeclareMathOperator{\Tor}{Tor}
\DeclareMathOperator{\PGL}{PGL}
\DeclareMathOperator{\hd}{hd}
\DeclareMathOperator{\pd}{pd}
\DeclareMathOperator{\SO}{SO}
\DeclareMathOperator{\SU}{SU}
\DeclareMathOperator{\Sp}{Sp}
\DeclareMathOperator{\Oth}{O}
\DeclareMathOperator{\Uni}{U}
\DeclareMathOperator{\evf}{ev}
\DeclareMathOperator{\cH}{H}
\DeclareMathOperator{\dR}{R}
\DeclareMathOperator{\End}{End}
\DeclareMathOperator{\Hasse}{Hasse}
\DeclareMathOperator{\Num}{Num}

\author{James Lee}

% Solarized
% \pagecolor[RGB]{0,20,26}
% \color[RGB]{191,191,191}

% Sephia
% \pagecolor[RGB]{249,239,220}

% Grey
% \pagecolor[RGB]{200, 200, 200}

% Black
% \pagecolor[RGB]{0,0,0}
% \color[RGB]{255,255,255}


\title{Chapter 5, Section 2}

\begin{document}

\maketitle

\begin{enumerate} [label=\textbf{\arabic*.}, leftmargin=0em]

  \item If $X$ is a birationally ruled surface, show that the curve $C$, such that $X$ is birationally equivalent to $C \times \P^1$, is unique up to isomorphism.

        \begin{proof}
          Let $C$ and $C'$ be two nonsingular curves such that $X$ is birational to $C \times \P^1$ and $C' \times \P^1$.
          Then, there exists rational maps $\pi, \pi' : X \dasharrow C, C'$ and sections $\sigma, \sigma' : C, C' \to X$ such that $\pi \circ \sigma$ and $\pi' \circ \sigma'$ extend to the identity morphism on $C$ and $C'$, respectively.
          Then $\pi' \circ \sigma : C \to C'$ and $\pi \circ \sigma' : C' \to C$ are inverses to each other, so we have an isomorphism $C \cong C'$.
        \end{proof}

  \item Let $X$ be the ruled surface $\P(\fE)$ over a curve $C$.
        Show that $\fE$ is decomposable if and only if there exist two sections $C', C''$ of $X$ such that $C' \cap C'' = \emptyset$.

        \begin{proof}
          In the forward direction, write $\fE = \fO \oplus \fO_C(\ge)$, where $e = - \deg{\ge} \leq 0$ (2.12).
          By (2.8.1), (2.12), let $C_0$ be a section corresponding to $\fE \to \fO_X(\ge)$, where $C_0^2 = -e$ (2.9).
          The surjective morphism of locally free sheaves $\fE \to \fO_C(\ge), \fO_C$ correspond to sections $C, C'$ linearly equivalent to $C_0$ and $C_0 - \ge f$ by (2.11.1), (2.11.2).
          Hence, the statement follows by $C_0.(C_0 - \ge f) = 0$.

          In the converse direction, let $C_0$ be a section as in (2.8.1), corresponding to a divisor $\ge$ of $C$ such that $\bigwedge^2 \fE \cong \fO_C(\ge)$, $e = - \deg{\ge}$, and $C_0^2 = -e$.
          By (2.9), there exists divisors $\gd', \gd''$ of $C$ corresponding to sections $C', C''$ of degree $d', d''$, respectively such that
          \begin{equation*}
            C' \sim C_0 + (\gd' - \ge)f, \quad C'' \sim C_0 + (\gd'' - \ge)f.
          \end{equation*}
          By (2.20), we have $d', d'' \geq 0$.
          If $C'$ and $C''$ do not intersect, then
          \begin{equation*}
            0 = C'.C'' = d' + d'' + e,
          \end{equation*}
          which shows $e \leq 0$.
          There are the surjecitve morphisms $\fE \to \fO_C(\gd'), \fO_C(\gd'')$
        \end{proof}

  \item \begin{itemize}
          \item[(a)] If $\fE$ is a locally free sheaf of rank $r$ on a nonsingular curve $C$, then there is a sequence
                \begin{equation*}
                  0 = \fE_0 \subseteq \fE_1 \subseteq \cdots \subseteq \fE_r = \fE
                \end{equation*}
                of subsheaves such that $\fE_i / \fE_{i - 1}$ is an invertible sheaf for each $i = 1, \dots, r$.
                We say that $\fE$ is a \textit{successive extension} of invertible sheaves.
          \item[(b)] Show that this is false for varieties of dimension $\geq 2$.
                In particular, the sheaf of differentials $\Omega_{\P^2}$ on $\P^2$ is not an extension of invertible sheaves.
        \end{itemize}

        \begin{proof} $ $ \vspace{0pt}
          \begin{itemize} [leftmargin=0cm]
            \item[(a)] We proceed by induction on $r$.
                  The base case $r = 1$ is trivial, so assume the statement for any $r > 1$.
                  We can twist $\fE$ by a very ample sheaf $\fL$ so that $\fE' = \fE \otimes \fE$ is generated by global sections.
                  If we prove the statement for $\fE'$ so we have a filtration $(\fE'_i)$, then tensoring by $\fL^{-1}$ gives a filtration $(\fE_i = \fE'_i \otimes \fL^{-1})$ of $\fE$ such that $$\fE_i / \fE_{i - 1} \cong (\fE_i' / \fE'_{i - 1}) \otimes \fL^{-1},$$
                  which is an invertible sheaf for each $i = 1, \dots, r$.
                  Hence, we can assume $\fE$ is generated by global sections.
                  The hypothesis of (II, Ex. 8.2) is satisfied for $\fE$, so there exists a short exact
                  \[ \begin{tikzcd}
                      0 \arrow[r] & \fO_C \arrow[r] & \fE \arrow[r] & \fE' \arrow[r] & 0,
                    \end{tikzcd} \]
                  where $\fE'$ is also locally free of rank $r - 1$.
                  By the inductive hypothesis, this means we can find a filtration $(\fE_i')$, $i = 0, \dots, r - 1$ of $\fE'$ such that each successive quotient is an invertible sheaf.
                  Consider the preimage $(\fE_j)$ of $(\fE_i')$, where we relabel $j = i + 1$.
                  Each $\fE_j$ is also locally free, $\fE_1 = \im(\fO_C)$ by exactness, and there is a canonical isomorphism $\fE_j / \fE_{j - 1} \cong \fE'_{j - 1} / \fE'_{j - 2}$.
                  Extending $(\fE_j)$ with $\fE_0 = 0$ gives a successive extension of sheaves.

            \item[(b)] Suppose we have an exact sequence
                  \[ \begin{tikzcd}
                      0 \arrow[r] & \fL \arrow[r] & \Omega \arrow[r] & \fM \arrow[r] & 0,
                    \end{tikzcd} \]
                  where $\fL$ and $\fM$ are invertible sheaves on $\P^2$.
                  Then $\cH^1(\fL), \cH^1(\fM) = 0$, so by the long exact of cohomology, $\cH^1(\Omega) = 0$.
                  By (III, Ex. 7.2), $\cH^1(\Omega) = k$, a contradiction.
          \end{itemize}
        \end{proof}

  \item Let $C$ be a curve of genus $g$, and let $X$ be the ruled surface $C \times \P^1$.
        We consider the question, for what integers $s \in \Z$ does there exist a section $D$ of $X$ with $D^2 = s$?
        First show that $s$ is always an even integer, say $s = 2r$.
        \begin{itemize}
          \item[(a)] Show that $r = 0$ and $r \geq g + 1$ are always possible.
          \item[(b)] If $g = 3$, show that $r = 1$ is not possible, and just one of the two values $r = 2, 3$ is possible, depending on whether $C$ is hyperelliptic or not.
        \end{itemize}

        \begin{proof}
          Fix a section $C_0 = C \times P_0 \subset X$, where $P_0$ is a point of $\P^1$.
          The ruled surface $X = C \times \P^1$ corresponds to the projective bundle of $\fE = \fO_C \oplus \fO_C$.
          The invariant of $X$ is $e = 0$, which means $C_0^2 = -e = 0$.
          By (2.9), if $D$ is any section of $X$ corresponding to a surjection $\fE \to \fO_C(\gd)$ for some divisor $\gd$ of $C$ with $r = \deg \gd$, then $r = C_0.D$ and $D \sim C_0 + \gd f$.
          Hence, $s = 2r$.

          \begin{itemize} [leftmargin=0cm]
            \item[(a)] There is a one-to-one correspondence between sections $C \to D \subset X$ and finite morphism $f : C \to \P^1$.
            In particular, a section $D \subset C \times \P^1$ is the graph of a morphism $f : C \to \P^1$, and its self-intersection is $D^2 = 2 \cdot \deg{f}$.
                  By (IV, Ex. 1.6), $g + 1$ is an upperbound on the minimum of $\deg{f}$.
                  Hence, $r = 0$ and $r \geq g + 1$ is always possible.

            \item[(b)] If $g = 3$, then there is no morphism $C \to \P^1$ of degree 1.
                  If $C$ is hyperelliptic, then it has a $g^1_2$, so $r = 2, 3$ is possible.
                  Otherwise, nonhyperelliptic curves of genus 3 always have a $g^1_3$ (IV, 5.5.2).
          \end{itemize}
        \end{proof}

        \newpage

  \item \textit{Values of $e$.} Let $C$ be a curve of genus $g \geq 1$.
        \begin{itemize}
          \item[(a)] Show that for each $0 \leq e \leq 2g - 2$ there is a ruled suface $X$ over $C$ with invariant $e$, corresponding to an indecomposable $\fE$.
          \item[(b)] Let $e < 0$, let $D$ be any divisor of degree $d = -e$, and let $\xi \in \cH^1(X, \fO_X(-D))$ be a nonzero element defining an extension
                \[ \begin{tikzcd}
                    0 \arrow[r] & \fO_C \arrow[r] & \fE \arrow[r] & \fO_X(D) \arrow[r] & 0.
                  \end{tikzcd} \]
                Let $H \subseteq |D + K|$ be the sublinear system of codimension 1 defined by $\ker{\xi}$, where $\xi$ is considered as a linear functional on $\cH^0(X, \fO_X(D + K))$.
                For any effective divisor $E$ of degree $d - 1$, let $L_E \subseteq |D + K|$ be the sublinear system $|D + K - E| + E$.
                Show that $\fE$ is normalized if and only if for each $E$ as above, $L_E \nsubseteq H$.
          \item[(c)] Now show that if $-g \leq e < 0$, there exists a ruled surface $X$ over $C$ with invariant $e$.
          \item[(d)] For $g = 2$, show that $e \geq -2$ is also necessary for the existence of $X$.
        \end{itemize}

        \begin{proof} $ $ \vspace{0pt}
          \begin{itemize} [leftmargin=0cm]
            \item[(a)]

            \item[(b)]

            \item[(c)]

            \item[(d)]
          \end{itemize}
        \end{proof}

  \item Show that every locally free sheaf of finite rank on $\P^1$ is isomorphic to a direct sum of invertible sheaves.

        \begin{proof}
          Let $\fE$ be a locally free sheaf of finite rank $r$ on $\P^1$.
          We proceed by induction on $r$.
          The case $r = 1$ is trivial.
          Assume the statement for some $r \geq 1$.
        \end{proof}

  \item On the elliptic ruled surface $X$ of (2.11.6), show that the sections $C_0$ with $C_0^2 = 1$ form a one-dimensional algebraic family, parametrized by the points of the base curve $C$, and that no two are linearly equivalent.

        \begin{proof}

        \end{proof}

  \item A locally free sheaf $\fE$ on a curve $C$ is said to be \textit{stable} if for every quotient locally free sheaf $\fE \to \fF \to 0$, $\fF \neq \fE$, $\fF \neq 0$, we have
        \begin{equation*}
          (\deg{\fF}) / \rank{\fF} > (\deg{\fE}) / \rank{\fE}
        \end{equation*}
        Replacing $>$ by $\geq$ defines \textit{semistable.}
        \begin{itemize}
          \item[(a)] A decomposable $\fE$ is never stable.
          \item[(b)] If $\fE$ has rank 2 and is normalized, then $\fE$ is stable (respectively, semistable) if and only if $\deg{\fE} . 0$ (respectively, $\geq 0$).
          \item[(c)] Show that the indecomposable locally free sheaves $\fE$ of rank 2 that are not seistable are classified, up to isomorphism, by giving (1) an integer $0 < e \leq 2g - 2$, (2) an element $\fL \in \pic{C}$ of degree $-e$, and (3) a nonzero $\xi \in \cH^1(X, \fL^{-1})$, deterimned up to a nonzero scalar multiple.
        \end{itemize}

        \newpage

  \item Let $Y$ be a nonsingular curve on a quadric cone $X_0$ in $\P^3$.
        Show that either $Y$ is a complete intersection of $X_0$ with a surface of degree $a \geq 1$, in which case $\deg{Y} = 2a$, $g(Y) = (a - 1)^2$, or $\deg{Y}$ is odd, say $2a + 1$, and $g(Y) = a^2 - a$.

  \item For any $n > e \geq 0$, let $X$ be the rational scroll of degree $d = 2n - e$ in $\P^{d + 1}$ given by (2.19).
        If $n \geq 2e - 2$, show that $X$ contains a nonsingular curve $Y$ of genus $g = d + 2$ which is a canonical curve in this embedding.
        Conclude that for every $g \geq 4$, there exists a nonhyperelliptic curve which has a $g_3^1$.

        \begin{proof}
          Fix a section $C_0 \subset X$ so that $C_0^2 = -e$.
          The embedding of $X$ in $\P^{d + 1}$ is given by the very ample divisor $D = C_0 + nf$.
          If $D' = aC_0 + bf$ is any divisor on $X$, then
          \begin{equation*}
            2p_a(D') - 2 = D'.(D' + K),
          \end{equation*}
          where $K = -2C_0 + (e - 2)f$ is a canonical divisor of $X$.
          Expanding the right-hand side and solving for $p_a(D')$ gives
          \begin{equation*}
            p_a(D') = \frac{1}{2}ea(3 - a) + (a - 1)(b - 1).
          \end{equation*}
          If $D'$ is irreducible and nonsingular, then its degree as a curve in $\P^{d + 1}$ equals to
          \begin{equation*}
            D.D' = (C_0 + nf).(aC_0 + bf) = (n - e)a + b.
          \end{equation*}
        \end{proof}

  \item Let $X$ be a ruled surface over the curve $C$, defined by a normalized bundle $\fE$, and let $\ge$ be the divisor on $C$ for which $\fO_C(\ge) \cong \bigwedge^2 \fE$ (2.8.1).
        Let $\gb$ be any divisor on $C$.
        \begin{itemize}
          \item[(a)] If $|\gb|$ and $|\gb + \ge|$ have no base points, and if $\gb$ is nonspecial, then there is a section $D \sim C_0 + \gb f$, and $|D|$ has no points.
          \item[(b)] If $\goth{b}$ and $\goth{b} + \ge$ are very ample on $C$, and for every point $P \in C$, we have $\gb - P$ and $\gb + \ge - P$ nonspecial, then $C_0 + \gb f$ is very ample.
        \end{itemize}

  \item Let $X$ be a ruled surface with invariant $e$ over an elliptic curve $C$, and let $\goth{b}$ be a divisor on $C$.
        \begin{itemize}
          \item[(a)] If $\deg{\gb} \geq \ge + 2$, then there is a section $D \sim C_0 + \goth{b}f$ such that $|D|$ has no base points.
          \item[(b)] The linear system $|C_0 + \gb f|$ is very ample if and only if $\deg{\gb} \geq e + 3$.
        \end{itemize}

  \item For every $e \geq -1$ and $n \geq e + 3$, there is an elliptic scroll of degree $d = 2n - e$ in $\P^{d _ 1}$.
        In particular, there is an elliptic scroll of degree 5 in $\P^4$.

  \item Let $X$ be a ruled surface over a curve of genus $g$, with invariant $e < 0$, and assume that $\car{k} = p > 0$ and $g \geq 2$.
        \begin{itemize}
          \item[(a)] If $Y \equiv aC_0 + bf$ is an irreducible curve $\neq C_0. f$, then either $a = 1$, $b \geq 0$, or $2 \leq a \leq p - 1$, $b \geq \frac{1}{2}ae$, or $a \geq p$, $b \geq \frac{1}{2} ae + 1 - g$.
          \item[(b)] If $a > 0$ and $b > a(\frac{1}{2}e + (1/p)(g - 1))$, then any divisor $D \equiv aC_0 + bf$ is ample.
                On the other hand, if $D$ is ample, then $a > 0$ and $b > \frac{1}{2}ae$.
        \end{itemize}

  \item \textit{Funny behaviour in characteristic $p$.} Let $C$ be the plane curve $x^3y + y^3z + z^3x = 0$ over a field $k$ of characteristic 3.
        \begin{itemize}
          \item[(a)] Show that the action of the $k$-linear Frobenius morphism $f$ on $\cH^1(C, \fO_C)$ is identically 0.
          \item[(b)] Fix a point $P \in C$, and show that there is a nonzero $\xi \in \cH^1(C, \fO_C(-P))$ such that $f^*\xi = 0$ in $\cH^1(C, \fO_C(-3P))$.
          \item[(c)] Now let $\fE$ be defined by $\xi$ as an extension
                \[ \begin{tikzcd}
                    0 \arrow[r] & \fO_C \arrow[r] & \fE \arrow[r] & \fO_C(P) \arrow[r] & 0,
                  \end{tikzcd} \]
                and let $X$ be the corresponding ruled surface over $C$.
                Show that $X$ contains a nonsingular curve $Y \equiv 3C_0 - 3f$, such that $\pi : Y \to C$ is purely inseparable.
                Show that the divisor $D = 2C_0$ satisfies the hypothesis of (2.21b), but is not ample.
        \end{itemize}

        \begin{proof} $ $ \vspace{0pt}
          \begin{itemize} [leftmargin=0cm]
            \item[(a)] Consider the short exact sequence
                  \
            \item[(b)]
            \item[(c)]
          \end{itemize}
        \end{proof}

  \item Let $C$ be a nonsingular affine curve.
        Show that two locally free sheaves $\fE, \fE'$ of the same rank are isomorphic if and only if their classes in the Grothendieck group $K(X)$ are the same.
        This is false for a projective curve.

        \begin{proof}
          Let $A$ be a Dedekind domain finitely generated over $k$ such that $C \cong \spec{A}$.
          The category of $\fO_C$-modules is equivalent to the category of $A$-modules.
          Under this identification, $K(X)$ is naturally isomorphic to the Grothendieck group $K_0(A)$ of flat finitely generated $A$-modules.
          The structure theorem of $A$-modules states $K_0(A) \cong \Z \oplus \cl(A)$, where if $M$ is any flat finitely generated $A$-module, then $M \cong A^{r - 1} \oplus I$, where if $K$ is the fraction field of $A$, then $r= \dim_K M \otimes_A K$ is the rank of $M$, and $I$ is some fractional ideal in $\cl(A)$.
          If $N \cong A^{s - 1} \oplus J$ is another $A$-module, then $M \cong N$ if and only if $r = s$ and $I \cong J$ in $\cl(A)$.
        \end{proof}

  \item \begin{itemize}
          \item[(a)] Let $\varphi : \P_k^1 \to \P_k^3$ be the 3-uple embedding.
                Let $\fI$ be the sheaf of ideals of the twisted cubic curve $C$ which is the image of $\varphi$.
                Then $\fI / \fI^2$ is a locally free sheaf of rank 2 on $C$, so $\varphi^*(\fI / \fI^2)$ is a locally free sheaf of rank 2 on $\P^1$.
                By (2.14), therefore, $\varphi^*(\fI/\fI^2) \cong \fO(l) \oplus \fO(m)$ for some $l, m \in \Z$.
                Determine $l$ and $m$.
          \item[(b)] Repeat part (a) for the embedding $\varphi : \P^1 \to \P^3$ given by $x_0 = t^4, x_1 = t^3u, x_2 = tu^3, x_3 = u^4$, whose image is a nonsingular rational quartic curve.
        \end{itemize}

        \begin{proof} $ $ \vspace{0pt}
          \begin{itemize} [leftmargin=0cm]
            \item[(a)] The canonical sheaf of $\P^3$ restricts to $\fO_{\P^1}(-12)$, and the canonical sheaf of $\P^1$ is $\fO_{\P^1}(-2)$, so $\deg \fI/\fI^2 = -10$.
            \item[(b)]
          \end{itemize}
        \end{proof}

\end{enumerate}

\end{document}
